\chapter{The Routing Engine and VLM Specialists}
\label{ch:routing_vlm}

Once a multi-page industrial document has been successfully processed through the Layout Detection phase, the resulting geometric bounding boxes must be interpreted. For generic text paragraphs, standard text primitives are structurally sufficient. However, tabular data within engineering datasheets exists across a vast spectrum of typological complexity. 

This chapter details the architecture of the Algorithmic Routing Engine developed in this thesis. By computing vector line densities and geometrical word-gap ratios, the router dynamically classifies tabular instances and dispatches them to distinct inferential systems: deterministic grid extractors, key-value specialists, or generative Vision-Language Models (VLMs). Furthermore, this section describes the multi-layered verification strategy consisting of constrained JSON generation and an automated Quality Assurance (QA) loop that triggers fallback extraction strategies upon failure.

\section{The Algorithmic Router}
\label{sec:router_logic}

Industrial tables generally manifest in three structural configurations: Ruled grids (explicitly boxed matrices), Key-Value pairs (implicit 2-column alignments), and Complex grids (multi-lane tables often embedding external graphical structures). Applying a massive generative Vision-Language Model to a perfectly ruled 50-row table incurs unacceptable latency and massive hallucination risk. Conversely, applying an implicit spatial algorithm to a 3-lane technical table fails catastrophically.

To resolve this, an internal sorting orchestrator, the \textit{TableTypeRouter}, computes mathematical heuristics based on the native topological primitives (words and drawing vectors) residing within a defined geometric bounding box. The orchestrator generates two primary continuous values: $S_{\text{ruled}}$ and $S_{\text{kv}}$.

\begin{figure}[htbp]
\centering
\begin{tikzpicture}[
    node distance=1cm and 1.5cm,
    box/.style={rectangle, draw, thick, align=center, minimum width=2.5cm, minimum height=0.8cm, rounded corners=0.1cm},
    decision/.style={diamond, draw, thick, align=center, aspect=2, minimum width=2.8cm, inner sep=0pt},
    arrow/.style={->, thick, >=stealth}
]

% Nodes
\node[box] (input) {Table Region $B_{table}$};
\node[box, below=0.6cm of input] (scores) {Compute $S_{ruled}$, $S_{kv}$};
\node[decision, below=0.8cm of scores] (is_ruled) {$S_{ruled} > 0.4$?};
\node[decision, below=1.2cm of is_ruled] (is_kv) {$S_{kv} > 0.8$?};

\node[box, right=1.5cm of is_ruled, fill=blue!10] (ruled) {Deterministic\\Grid Extractor};
\node[box, right=1.5cm of is_kv, fill=green!10] (vlm) {Specialist VLM\\(JSON Mode)};

\node[box, below=1.2cm of is_kv, fill=red!10] (tsr) {Complex TSR\\\& Masking};

\node[box, below=1.0cm of vlm, fill=orange!10] (qa) {QA Layer\\(Coverage/Duplication)};

% Connections
\draw[arrow] (input) -- (scores);
\draw[arrow] (scores) -- (is_ruled);

\draw[arrow] (is_ruled) -- node[above, font=\small] {Yes} (ruled);
\draw[arrow] (is_ruled) -- node[left, font=\small] {No} (is_kv);

\draw[arrow] (is_kv) -- node[above, font=\small] {Yes} (vlm);
\draw[arrow] (is_kv) -- node[left, font=\small] {No} (tsr);

\draw[arrow] (ruled) |- (qa);
\draw[arrow] (vlm) -- (qa);

\draw[arrow] (qa) -- node[right, font=\small] {Fail $\rightarrow$ Retry} +(2,0) |- (scores);

\end{tikzpicture}
\caption{The hierarchical routing and verification loop. Tables are routed based on geometric scores; the extraction is subsequently validated by a QA layer that can trigger fallback attempts.}
\label{fig:router_flow}
\end{figure}

\subsection{Vector Line Density Scoring}
To determine if a tabular bounding box contains a discrete physical grid, the router calculates the cumulative Euclidean length of all valid horizontal and vertical drawing primitives ($D$). Let $B$ represent the table bounding box with width $w$ and height $h$. The perimeter is defined as $P = 2(w + h)$. The ruled score $S_{\text{ruled}}$ is normalized against the bounded perimeter:

\begin{equation}
S_{\text{ruled}} = \min \left( 1.0, \frac{\sum_{d \in D} \text{length}(d)}{2(w + h)} \right)
\end{equation}

If $S_{\text{ruled}} > R_{threshold}$, the bounding box is immediately classified as a physically ruled table, triggering deterministic row-column projection parsing.

\subsection{Maximum Gap Ratio Scoring}
For technical datasheets heavily utilizing implicit Key-Value mappings (e.g., "Operating Voltage | 24 V DC"), physical lines are entirely absent. To detect this 2-column configuration without relying on an LLM, the router analyzes the geometric distribution of the textual nodes.

Let $X$ represent the sorted array of the geometric $X$-centers corresponding to every word bounding box within the table. The continuous gaps $G$ between adjacent words are computed. The maximum topological gap $g_{max}$ and its positional ratio relative to the table width $w$ form the basis of the key-value score $S_{\text{kv}}$:

\begin{equation}
S_{\text{kv}} = 0.7 \left( \frac{g_{max}}{w} \right) + 0.3 \left( B_{factor} \right)
\end{equation}

The balance factor, $B_{factor}$, calculates the symmetric distribution of text nodes residing left or right of $g_{max}$. If a table yields a highly imbalanced distribution or exhibits multiple significant gaps ($n \geq 2$), the router detects a complex multi-metric grid, bypassing the VLM and forcing the table to the Semantic Pre-Masking TSR framework detailed in Chapter 5.

\section{Structured Extraction via JSON Schemas}
\label{sec:vlm_specialist}

When the router encounters a structure that requires generative reasoning ($S_{\text{kv}} > K_{threshold}$), the region is passed to a Vision-Language Extractor. Unlike general-purpose chat models, the extractor used in this pipeline is constrained to produce output that adheres monotonically to a predefined Pydantic schema.

\subsection{Constrained Generation and JSON Mode}
To ensure the output is machine-readable and semantically valid, the pipeline utilizes "JSON Mode" combined with strict schema definitions \cite{json_mode}. The extractor receives a system prompt containing a technical taxonomy (e.g., the \textit{LibrarianUniversalHardware} schema) and is instructed to populate only the fields identified within the document.

This approach offers two advantages over legacy XML-tagging. First, it enables complex nested structures, such as pin-to-signal assignments for industrial connectors, which are difficult to parse with regular expressions in XML formats. Second, it allows for type-level validation; the pipeline can verify that "Supply Voltage" is a numeric range rather than an arbitrary string at the point of inference.

\section{Extraction Quality Assurance and Adaptive Retry}
\label{sec:qa_layer}

A significant contribution of this architecture is the secondary verification layer that follows every extraction attempt. Rather than assuming the first extraction is correct, the \textit{QA Coordinator} evaluates two precision metrics: \textbf{Coverage} and \textbf{Duplication}.

\subsection{Coverage and Duplication Metrics}
The coverage metric measures the ratio of source text tokens successfully mapped into the resulting JSON structure. If a table contains forty numeric parameters but the VLM only extracts ten, the low coverage score triggers an automated failure.

Conversely, the duplication metric detects instances where the model "stutters" or recursively generates the same row multiple times --- a common failure mode in sequence-to-sequence document models. Let $E$ be the set of extracted records. The duplication score $S_{dup}$ is defined as the frequency of redundant entries normalized by the total count. If $S_{dup} > 0.05$ or Coverage $< 0.90$, the extraction is rejected.

\subsection{Adaptive Fallback Logic}
In the event of a QA failure, the pipeline does not stop. Instead, the coordinator triggers an \textit{Adaptive Fallback}. If a generative VLM attempt failed due to duplication error, the system may re-route the region to a deterministic grid parser or attempt the extraction again with a higher "temperature" or a modified prompt structure that emphasizes structural unique-constraints. This iterative retry loop substantially increases the extraction success rate on highly complex, semi-structured industrial materials.

\subsection{Eradicating Generative Topological Hallucination}
The most pernicious failure state in visual grid reconstruction occurs when a standard VLM processes a technical table encapsulating a visual diagram (e.g., an embedded sensor wiring schematic). Standard instruction-tuned models recursively attempt to OCR the diagram's internal axes or numeric labels, creating endlessly hallucinated, unsynchronized tabular rows. 

To surgically mitigate this hallucination loop, the XML constraint strategy is augmented with a categorical override variable. The VLM inference generation prompt defines a critical exit condition: \textit{If a row contains a FIGURE or DIAGRAM, absolutely ignore all internal text and structural coordinates. Output \texttt{<V>[TECHNICAL\_GRAPH]</V>} and cease iteration.} 

This semantic intervention proves phenomenally successful. Empirical evaluation reveals that providing the Generative model with an explicit semantic exit-tag dramatically suppresses its latent drive to infinitely parse arbitrary graphical lines, effectively stabilizing the 2-column key-value extraction pipeline.
